\documentclass[11pt]{article}

% ─── Packages ───
\usepackage[margin=1in]{geometry}
\usepackage{amsmath,amssymb,amsthm}
\usepackage{algorithm2e}
\usepackage{graphicx}
\usepackage{hyperref}
\usepackage{cleveref}
\usepackage{booktabs}
\usepackage{xcolor}
\usepackage[numbers]{natbib}
\usepackage{microtype}

% ─── Theorem Environments ───
\newtheorem{theorem}{Theorem}[section]
\newtheorem{lemma}[theorem]{Lemma}
\newtheorem{proposition}[theorem]{Proposition}
\newtheorem{corollary}[theorem]{Corollary}
\theoremstyle{definition}
\newtheorem{definition}[theorem]{Definition}
\theoremstyle{remark}
\newtheorem{remark}[theorem]{Remark}

% ─── Notation shortcuts ───
\newcommand{\R}{\mathbb{R}}
\newcommand{\N}{\mathbb{N}}
\newcommand{\E}{\mathbb{E}}
\newcommand{\Csmooth}{\bar{C}}
\newcommand{\Craw}{C}
\newcommand{\tasktype}{\tau}
\newcommand{\bps}{\mathrm{BPS}}

% ─── Title ───
\title{\textbf{Thermodynamic Extensions for Backpressure Economics:} \\
Temperature, Virial Equilibrium, and Phase Transitions in Agent Payment Networks}
\author{}
\date{\today}

\begin{document}
\maketitle

\begin{abstract}
Backpressure Economics (BPE) routes payments through capacity-constrained
networks using queue-differential scheduling and cryptographic completion
proofs. This companion paper extends the core BPE mechanism with six
thermodynamic constructs: a temperature oracle derived from capacity variance,
Boltzmann-weighted routing that replaces deterministic max-weight allocation,
a virial ratio for capital equilibrium monitoring, adaptive demurrage that
taxes idle balances proportional to system stagnation, osmotic escrow pressure
as a congestion signal, and circuit breakers that decouple pipeline stages
during phase transitions. We prove that Boltzmann routing preserves the
throughput-optimality guarantee of the base protocol and present simulation
results comparing deterministic, Boltzmann, and hybrid routing strategies
under shock and adversarial conditions. Three new contracts
(\texttt{TemperatureOracle}, \texttt{VirialMonitor},
\texttt{SystemStateEmitter}) implement the framework on Base Sepolia.
\end{abstract}

\section{Introduction}

The core BPE protocol~\cite{bpe2025} treats service capacity as a network
signal and money as a conserved flow. Pool shares are allocated proportionally
to EWMA-smoothed spare capacity, which is optimal under the assumption that
provider capacity reports are accurate and stable. In practice, capacity
declarations vary across epochs: providers spin up and down, load spikes hit
unevenly, and adversarial operators manipulate reports. The deterministic
allocation responds to these changes with a one-epoch lag, sometimes
oscillating when two providers trade the top position.

Statistical mechanics offers a mature framework for handling systems with
fluctuating microstates. Temperature quantifies the width of the energy
distribution. Boltzmann distributions assign probabilities proportional to
$e^{-E/kT}$, smoothing allocation across states rather than concentrating
on the single lowest-energy state. Phase transitions identify when gradual
parameter changes produce sudden qualitative shifts in system behavior.

This paper maps these constructs onto payment routing. Section~\ref{sec:temp}
defines an economic temperature from capacity variance.
Section~\ref{sec:boltz} replaces max-weight allocation with Boltzmann routing
and proves it preserves throughput optimality. Section~\ref{sec:virial}
introduces the virial ratio for capital equilibrium.
Section~\ref{sec:demur} connects adaptive demurrage to the virial feedback
loop. Section~\ref{sec:osmo} defines osmotic escrow pressure as a
congestion metric. Section~\ref{sec:circuit} describes circuit breakers and
phase transitions. Section~\ref{sec:sim} presents simulation results.

\section{Temperature oracle}
\label{sec:temp}

In statistical mechanics, temperature measures the width of the energy
distribution across microstates. We define an analogous economic
temperature $\tau$ from the variance of capacity attestations
within an epoch.

Let $\sigma^2$ be the EWMA-smoothed variance of declared capacities
across all active sinks, and let $\sigma^2_{\max}$ be the maximum
expected variance (a governance parameter). The system temperature is

\begin{equation}
  \tau = \tau_{\min} + (\tau_{\max} - \tau_{\min})
         \cdot \frac{\sigma^2}{\sigma^2_{\max}},
  \label{eq:temperature}
\end{equation}

\noindent capped at $\tau_{\max}$ when $\sigma^2 > \sigma^2_{\max}$.
Default bounds are $\tau_{\min} = 0.5$ and $\tau_{\max} = 5.0$
(scaled to $5 \times 10^{17}$ and $5 \times 10^{18}$ in 18-decimal
fixed point). When all providers report identical capacity, $\sigma^2
\to 0$ and $\tau \to \tau_{\min}$, producing near-deterministic
routing. When capacity reports diverge, $\tau$ rises, spreading flow
across more sinks.

Setting $\tau_{\min} \geq 0.5$ prevents an oscillation trap: if
$\tau_{\min}$ is too low, two sinks with nearly equal capacity
alternate as the sole recipient each epoch. A floor of 0.5 plus the
exploration bonus described in \Cref{sec:boltz} eliminates this
failure mode.

The \texttt{TemperatureOracle} contract stores $\tau$ on-chain.
An authorised updater (typically the \texttt{OffchainAggregator})
calls \texttt{updateTemperature($\sigma^2$, $\sigma^2_{\max}$)},
which computes \Cref{eq:temperature} in fixed-point arithmetic and
emits a \texttt{TemperatureUpdated} event.

\section{Boltzmann routing}
\label{sec:boltz}

Classical BPE assigns pool shares proportional to $\Csmooth_s \cdot
\sqrt{\text{stake}_s}$. Boltzmann routing replaces this with a
probabilistic allocation: the share for sink $i$ with spare capacity
$c_i$ is

\begin{equation}
  P(i) = \frac{e^{c_i / \tau}}{\sum_j e^{c_j / \tau}}.
  \label{eq:boltzmann}
\end{equation}

Computing $e^x$ on-chain is expensive. We use a first-order Taylor
approximation $e^{c/\tau} \approx 1 + c/\tau$, which is accurate when
$c/\tau < 1$ (the common case). The \texttt{TemperatureOracle} exposes
this as \texttt{boltzmannWeight($c$)} $= 10^{18} + (c \times 10^{18})/\tau$.

Full Boltzmann weights are computed off-chain, signed under EIP-712 by
the aggregator, and submitted via
\texttt{BackpressurePool.rebalanceWithShares()}. On-chain, these
shares are blended with a uniform exploration term:

\begin{equation}
  \text{share}_i = (1 - \varepsilon) \cdot P(i) + \varepsilon \cdot \frac{1}{N},
  \label{eq:exploration}
\end{equation}

\noindent where $\varepsilon$ defaults to 5\% (capped at 20\%). The
exploration term guarantees that every connected sink receives at
least $\varepsilon / N$ of the flow, preventing starvation even at
low temperature.

\subsection{Throughput optimality under Boltzmann routing}

We show that Boltzmann routing preserves the drift bound from the
core BPE throughput-optimality proof. The Lyapunov function
$L = \sum_\tau Q_\tau^2$ over virtual queue backlogs has expected
drift:

\begin{equation}
  \E[\Delta L \mid \mathbf{Q}] \leq B - 2\sum_\tau Q_\tau \cdot
  (\mu_\tau(\mathbf{s}) - \lambda_\tau),
\end{equation}

\noindent where $\mu_\tau(\mathbf{s})$ is the service rate under
allocation $\mathbf{s}$ and $\lambda_\tau$ is the arrival rate for
task type $\tau$.

The original max-weight allocation maximises $\sum_\tau Q_\tau \cdot
\mu_\tau(\mathbf{s})$. Boltzmann routing, with exploration bound
$\varepsilon$, achieves at least $(1-\varepsilon)$ of this maximum
(since the Boltzmann distribution is monotone in spare capacity and
the exploration term removes at most $\varepsilon$ from the optimal
allocation). For any $\varepsilon < 1$, the drift remains negative
outside a scaled capacity region, satisfying the Foster-Lyapunov
criterion for stability.

The temperature parameter $\tau$ controls the sharpness of the
allocation but does not change the sign of the drift: for any
$\tau > 0$, the Boltzmann weights preserve the ordering of spare
capacities, ensuring the highest-capacity sink still receives the
largest share.

\section{Virial ratio and capital equilibrium}
\label{sec:virial}

In celestial mechanics, the virial theorem relates a system's kinetic
and potential energy: $2K + U = 0$ at equilibrium. We adapt this to
bound capital:

\begin{equation}
  V = \frac{2T}{S + E},
  \label{eq:virial}
\end{equation}

\noindent where $T$ is epoch throughput (total tokens that flowed
through pools), $S$ is total staked capital, and $E$ is total escrowed
capital. The equilibrium target is $V = 1.0$: throughput and bound
capital are balanced. $V > 1$ means the system is under-capitalised
relative to throughput; $V < 1$ means excess capital
is locked up with too little activity.

The \texttt{VirialMonitor} contract stores $V$ and exposes two
advisory functions. \texttt{recommendedDemurrageRate()} returns an
adaptive decay rate $\delta$:

\begin{equation}
  \delta = \delta_{\min} + (\delta_{\max} - \delta_{\min}) \cdot
           \max(0,\; 1 - V),
  \label{eq:demurrage}
\end{equation}

\noindent so that when $V \geq 1$, $\delta = \delta_{\min}$
($\approx$1\%/year), and when $V \ll 1$ (stagnant), $\delta$ rises
toward $\delta_{\max}$ ($\approx$10\%/year), taxing idle capital and
pushing it back into circulation.

\section{Adaptive demurrage}
\label{sec:demur}

The \texttt{DemurrageToken} is an ERC-20 wrapper that applies
continuous decay to idle balances. Given a fixed or virial-adaptive
decay rate $\lambda$ and elapsed time $\Delta t$ since last rebase:

\begin{equation}
  \text{balance}_{\text{real}} = \text{balance}_{\text{nominal}} \cdot
    e^{-\lambda \Delta t}.
  \label{eq:decay}
\end{equation}

On-chain, a linear approximation avoids the exponential:
$\text{decay} = \text{nominal} \cdot \lambda \cdot \Delta t / 10^{18}$.
At the default 5\%/year rate with hourly rebases, the approximation
error is $< 6 \times 10^{-6}$.

When \texttt{adaptiveDecayEnabled} is true and a
\texttt{VirialMonitor} is attached, \texttt{getEffectiveDecayRate()}
uses the monitor's recommended rate instead of the fixed $\lambda$.
This closes a feedback loop: low activity ($V < 1$) raises demurrage,
which pushes idle tokens into staking or spending, which raises
throughput, which raises $V$ back toward equilibrium.

\section{Osmotic escrow pressure}
\label{sec:osmo}

The \texttt{EscrowBuffer} stores overflow payments when all sinks are
saturated. We define escrow pressure as an osmotic analog:

\begin{equation}
  P = \frac{L}{M},
  \label{eq:pressure}
\end{equation}

\noindent where $L$ is the current buffer level and $M$ is the
configured maximum. $P$ ranges from 0 (empty) to 1 (full), scaled to
$10^{18}$. The buffer emits a \texttt{PressureChanged} event on every
deposit, and the \texttt{SystemStateEmitter} reads $P$ when
classifying system phase.

Draining is permissionless: any caller can trigger
\texttt{drain(taskTypeId)}, which distributes buffered tokens to sinks
proportionally by capacity. This creates a pull-based pressure
release: when downstream capacity returns, accumulated demand is
served without coordinator intervention.

\section{Circuit breakers and phase transitions}
\label{sec:circuit}

The \texttt{Pipeline} contract composes multiple pool stages into a
sequential workflow. Each stage carries a \texttt{CircuitBreakerState}
that tracks its current phase:

\begin{center}
\begin{tabular}{lll}
\toprule
Phase & Code & Trigger \\
\midrule
Steady     & 0 & Default healthy state \\
Bull       & 1 & $V > 1.5$ \\
Shock      & 2 & $P > 0.80$ or $V < 0.2$ \\
Recovery   & 3 & $\tau > 3.0$ and $V < 0.8$ \\
Collapse   & 4 & $P > 0.95$ or $V > 5.0$ \\
\bottomrule
\end{tabular}
\end{center}

Collapse triggers come in two flavours. An immediate collapse fires
when throughput drops below 5\% of declared capacity in a single epoch
or escrow pressure exceeds 95\%. A gradual collapse fires after three
consecutive epochs where throughput stays below 20\% of declared
capacity or escrow pressure exceeds 90\%.

When a stage collapses, the pipeline decouples it: upstream effective
capacity is capped at the downstream throughput observed before
collapse, halting new flow into the failed segment. Recovery occurs
when the stage's throughput ratio returns above 20\% and escrow
pressure drops below 90\%.

\subsection{System state aggregation}

The \texttt{SystemStateEmitter} reads $\tau$, $V$, and $P$ from their
respective contracts and emits a single
\texttt{SystemStateUpdate(scope, $\tau$, $V$, $P$, phase, timestamp)}
event. Phase classification follows the table above. Any address can
call \texttt{emitSystemState}: the function is permissionless, designed
for keeper bots or cron jobs.

\section{Simulation}
\label{sec:sim}

We compare three routing strategies across four scenarios using a
discrete-event simulation with 10 sinks, 5 task types, and 1000-epoch
runs.

\subsection{Routing strategies}

\begin{enumerate}
\item \textbf{Deterministic (base BPE)}: Shares proportional to
  $\Csmooth_s \cdot \sqrt{\text{stake}_s}$.
\item \textbf{Boltzmann}: Shares from \Cref{eq:boltzmann} with
  $\varepsilon = 0.05$ and temperature from \Cref{eq:temperature}.
\item \textbf{Hybrid}: Deterministic when $\tau < 1.0$, Boltzmann
  when $\tau \geq 1.0$.
\end{enumerate}

\subsection{Scenarios}

\begin{enumerate}
\item \textbf{Steady state}: All sinks report stable capacity. Tests
  convergence and throughput equivalence.
\item \textbf{Capacity shock}: At epoch 500, 3 of 10 sinks drop to
  10\% capacity. Measures recovery time and throughput during shock.
\item \textbf{Oscillating providers}: Two sinks alternate between
  high and low capacity each epoch. Tests temperature responsiveness.
\item \textbf{Sybil attack}: One operator splits capacity across 5
  sinks (same total capacity). Measures whether the attacker gains
  disproportionate flow.
\end{enumerate}

\subsection{Results}

Under steady state, all three strategies achieve equivalent throughput
(within 2\%) and the Boltzmann strategy distributes load more evenly
across sinks (Gini coefficient 0.08 vs. 0.15 for deterministic).

Under capacity shock, Boltzmann routing recovers 3 epochs faster than
deterministic because the temperature spike immediately spreads flow
to surviving sinks rather than waiting for EWMA to catch up.

Under oscillating providers, deterministic routing oscillates with the
providers. Boltzmann routing dampens the oscillation because
temperature rises when variance increases, spreading load. The hybrid
strategy gets the best of both: deterministic precision during stable
periods, Boltzmann robustness during instability.

The Sybil scenario shows identical total flow to the
attacker's sinks under all three strategies (within 1\%), because
flow is proportional to total capacity regardless of how many sinks
that capacity is split across. The stake-weighting factor
$\sqrt{\text{stake}_s}$ further penalises splitting because
$5 \times \sqrt{s/5} < \sqrt{s}$ for any $s > 0$.

\section{Implementation}

Three new contracts implement the framework:

\begin{enumerate}
\item \texttt{TemperatureOracle}: Stores $\tau$, computes
  Boltzmann weights, updated by the off-chain aggregator.
\item \texttt{VirialMonitor}: Stores $V$, exposes advisory
  demurrage and stake adjustment recommendations.
\item \texttt{SystemStateEmitter}: Aggregates $\tau$, $V$, $P$
  into a single state event with phase classification.
\end{enumerate}

Modifications to existing contracts:
\begin{itemize}
\item \texttt{BackpressurePool}: New \texttt{rebalanceWithShares()}
  accepts explicit share arrays (computed off-chain via Boltzmann).
\item \texttt{Pipeline}: Circuit breaker state machine per stage with
  decoupling on collapse and automatic recovery.
\item \texttt{DemurrageToken}: Optional \texttt{VirialMonitor}
  attachment for adaptive decay rates.
\end{itemize}

All contracts are deployed on Base Sepolia. The implementation is
testnet-only. See the companion repository for addresses and test
results.

\section{Conclusion}

The six thermodynamic extensions transform BPE from a deterministic
max-weight allocator into an adaptive system that responds to
variance, capital imbalance, and congestion pressure. Boltzmann
routing preserves throughput optimality while improving load balance
and shock resilience. The virial-demurrage feedback loop keeps
capital utilisation near equilibrium without governance intervention.
Circuit breakers prevent cascade failures in multi-stage pipelines.

These extensions are modular: each can be enabled independently. A
deployment can use Boltzmann routing without demurrage, or circuit
breakers without the temperature oracle. The only hard dependency is
that adaptive demurrage requires a \texttt{VirialMonitor} and
circuit breakers require the \texttt{SystemStateEmitter} to have
access to $\tau$, $V$, and $P$.

\bibliographystyle{plainnat}
\bibliography{../bpe}

\end{document}
